\section{Conversion des données du contrôleur M221 (HMI)}
\label{secConversionHmi}

Le contrôleur M221 produit 4 octets et en consomme 40 (cf. document
\texttt{02b\_donnees\_echangees}, section « Données de l'IHM et du
pupitre opérateur »). Ces octets ne sont pas directement exploitables :
ils doivent être lus et écrits à travers les blocs fonctionnels maison
\texttt{FB\_ReadHmi} et \texttt{FB\_WriteHmi}.

\subsection{En lecture}

\begin{enumerate}
    \item L'alias \texttt{fromHmiScara}, de type \texttt{TFromHmi},
    chevauche la zone image des données produites par le M221 dans
    l'espace \texttt{\%MW} du coupleur \texttt{NOC\_ETHIP\_ROBOT}. Il ne
    contient que des octets bruts, sans nom de champ exploitable.
    \item Le bloc \texttt{FB\_ReadHmi} prend cet objet brut en entrée et
    en extrait les informations individuelles (chaque sélecteur, chaque
    bouton-poussoir) vers un objet \texttt{THmiInputs} dont les champs
    portent un nom explicite : \texttt{m\_nixJ1Plus}, \texttt{m\_nixAuto},
    \texttt{m\_nixStopRequest}...
\end{enumerate}

\texttt{FB\_ReadHmi} expose les entrées/sorties suivantes :

\begin{center}
\begin{tabular}{|l|l|p{7.3cm}|}
\hline
\textbf{Paramètre} & \textbf{Sens / type} & \textbf{Rôle} \\
\hline
\texttt{Enable} & \texttt{VAR\_INPUT}, \texttt{BOOL} & Active la
conversion à ce cycle (en général, toujours \texttt{TRUE}). \\
\hline
\texttt{IN} & \texttt{VAR\_INPUT}, \texttt{TFromHmi} & Objet brut,
chevauchant la zone \texttt{\%MW} image des données reçues du M221. \\
\hline
\texttt{Inputs} & \texttt{VAR\_OUTPUT}, \texttt{THmiInputs} & Objet
sémantique mis à jour par le bloc, exploitable par le reste du programme. \\
\hline
\texttt{Error} & \texttt{VAR\_OUTPUT}, \texttt{BOOL} & Signale une
erreur de conversion (ex. taille de \texttt{IN} incohérente). \\
\hline
\end{tabular}
\end{center}

\textbf{Exemple d'appel} (section exécutée en tout début de cycle) :

\begin{lstlisting}
InstReadHmi(
    Enable := TRUE,
    IN     := fromHmiScara,
    Inputs => hmiInputsScara,
    Error  => errorReadHmi
);
\end{lstlisting}

\begin{UPSTIinfor}{Pourquoi un type \texttt{THmiInputs} dédié ?}
\texttt{FB\_ReadHmi} pourrait, en théorie, exposer chaque bit décodé comme
une sortie \texttt{VAR\_OUTPUT} séparée (\texttt{BOOL} par
\texttt{BOOL}). Regrouper ces champs dans un \textbf{type structuré}
dédié, \texttt{THmiInputs}, présente plusieurs avantages concrets :
\begin{itemize}
    \item un \textbf{seul} paramètre suffit pour transmettre l'ensemble
    des informations à une autre partie du programme (une fonction qui
    gère le mode automatique n'a besoin que d'un argument
    \texttt{THmiInputs}, pas d'une dizaine de \texttt{BOOL} séparés) ;
    \item ajouter un nouveau champ (un futur bouton-poussoir) ne change ni
    la signature du bloc ni celle des fonctions qui consomment
    \texttt{THmiInputs} : seul le corps de \texttt{FB\_ReadHmi} est
    modifié ;
    \item le type documente, par ses noms de champs, la signification de
    chaque information -- contrairement à un tableau de bits anonyme.
\end{itemize}
\end{UPSTIinfor}

\begin{UPSTIinfor}{Attention au sens des contacts}
Certains champs de \texttt{THmiInputs} correspondent à des contacts à
ouverture (ex. \texttt{m\_nixStopRequest}, \texttt{m\_nixEmptyTank0}) :
une valeur \texttt{FALSE} y signifie que la condition surveillée est
\textbf{active}. \texttt{FB\_ReadHmi} ne fait qu'extraire le bit brut, il
n'inverse pas son sens -- c'est au programme applicatif d'en tenir
compte.
\end{UPSTIinfor}

\subsection{En écriture}

Le principe est symétrique : le programme applicatif renseigne les champs
de \texttt{THmiOutputs} (\texttt{m\_qToR} pour les voyants du pupitre,
\texttt{m\_Screen} pour la supervision écran -- \texttt{m\_iStateGmma},
\texttt{m\_iCurvSpeed\_mm\_per\_s}, \texttt{m\_a5tankLine}...), puis
\texttt{FB\_WriteHmi} se charge de reconstituer, à partir de ces champs,
l'objet brut \texttt{TToHmi} qui sera réellement écrit par le coupleur
vers le M221.

\begin{center}
\begin{tabular}{|l|l|p{7.3cm}|}
\hline
\textbf{Paramètre} & \textbf{Sens / type} & \textbf{Rôle} \\
\hline
\texttt{Enable} & \texttt{VAR\_INPUT}, \texttt{BOOL} & Active la
conversion à ce cycle. \\
\hline
\texttt{Outputs} & \texttt{VAR\_INPUT}, \texttt{THmiOutputs} & Objet
sémantique renseigné par le programme applicatif. \\
\hline
\texttt{OUT} & \texttt{VAR\_OUTPUT}, \texttt{TToHmi} & Objet brut produit
par le bloc, à écrire par le coupleur vers le M221. \\
\hline
\texttt{Error} & \texttt{VAR\_OUTPUT}, \texttt{BOOL} & Signale une
erreur de conversion. \\
\hline
\end{tabular}
\end{center}

\textbf{Exemple d'appel} (section exécutée en fin de cycle) :

\begin{lstlisting}
InstWriteHmi(
    Enable  := TRUE,
    Outputs := hmiOutputsScara,
    OUT     => toHmiScara,
    Error   => errorWriteHmi
);
\end{lstlisting}

\begin{UPSTIidee}{Trois temps, toujours dans le même ordre}
On retrouve, dans l'ordre, les trois étapes suivantes :
\begin{enumerate}
    \item \textbf{Lire} l'objet brut (alias sur la zone \texttt{\%MW})
    puis le \textbf{convertir} en objet sémantique (\texttt{FB\_ReadHmi}).
    \item \textbf{Utiliser} (ou renseigner) l'objet sémantique dans la
    logique applicative.
    \item \textbf{Convertir} l'objet sémantique en objet brut
    (\texttt{FB\_WriteHmi}) puis laisser le coupleur l'\textbf{écrire}
    physiquement.
\end{enumerate}
\end{UPSTIidee}

\UPSTIremarque[Lien avec la préparation du TP]{
C'est précisément ce que demande l'activité « Déclaration des objets
\texttt{TFromHmi} et \texttt{TToHmi} » du TP suivant : déclarer, par
chevauchement/alias, les objets \texttt{fromHmiScara} et
\texttt{toHmiScara} qui serviront d'entrée/sortie brutes à
\texttt{FB\_ReadHmi} et \texttt{FB\_WriteHmi}.}
